\documentclass[12pt]{article}                                                                                                           
\usepackage[utf8]{inputenc}                                                                                                             
\usepackage[T1]{fontenc}                                                                                                                
\usepackage{amsmath, amssymb}                                                                                                           
\usepackage{geometry}                                                                                                                   
\usepackage{xcolor}                                                                                                                     
\usepackage{lipsum}                                                                                                                     
\usepackage{titlesec}                                                                                                                   
\usepackage{fancyhdr}                                                                                                                   
\usepackage[most]{tcolorbox}                                                                                                                  
\usepackage{graphicx}                                                                                                                   
\usepackage{float}                                                                                                      
\usepackage{tikz}
\usetikzlibrary{shapes.geometric, arrows.meta}
\usetikzlibrary{arrows.meta, decorations.markings}                                                                                      
                                                                                                                                        
\geometry{a4paper, margin=2.5cm}                                                                                                        
\pagestyle{fancy}                                                                                                                       
\fancyhf{}                                                                                                                              
\rhead{Problem Set I}                                                                                                                   
\lhead{Rigid body mechanics}                                                                                                                     
\cfoot{\thepage}                                                                                                                        
                                                                                                                                        
% Fancy titles                                                                                                                          
\titleformat{\section}{\normalfont\Large\bfseries}{Exercise \thesection:}{1em}{}                                                        
\titleformat{\subsection}{\normalfont\bfseries}{Correction:}{1em}{}                                                                     
                                                                                                                                        
% Custom box for correction with breakable option
\newtcolorbox{correctionbox}{                                                                                                           
  colback=gray!5,                                                                                                                       
  colframe=black,                                                                                                                       
  fonttitle=\bfseries,                                                                                                                  
  title=Correction,                                                                                                                     
  before skip=10pt,                                                                                                                     
  after skip=10pt,
  breakable,
  enhanced
}                                                                                                                                       
                                                                                                                                        
% Fix headheight warning
\setlength{\headheight}{14.49998pt}

\begin{document}

\begin{center}
\Large\textbf{Ibn Tofail University} \\[1em]
\large\textit{Rigid body mechanics} \\[2em]
\large\textit{Problem Set I} \\[0.5em]
\end{center}

\vspace{1cm}
% ----------- EXERCISE 1 -------------
\section{}

Let $(A, \vec{V})$ be a bound vector defined in the basis $(\vec{i}, \vec{j}, \vec{k})$ by: $A(-2, -1, 2)$ and $\vec{V} = \vec{i} + 2\vec{j} - \vec{k}$.

\begin{enumerate}
    \item Calculate the moment of the vector $\vec{V}$ at the point with coordinates $(2, 1, 1)$.
    \item Calculate the moment of $(A, \vec{V})$ with respect to the axis $\Delta$ passing through $B$ and parallel to the vector $\vec{U} = 2\vec{i} - 3\vec{k}$.
\end{enumerate}

\begin{correctionbox}

\end{correctionbox}

% ----------- EXERCISE 2 -------------
\section{}

We are given two torsors $T_1$ and $T_2$. Find the set of points (if it exists) in space where their moments are collinear.

\begin{correctionbox}

\end{correctionbox}

% ----------- EXERCISE 3 -------------
\section{}

Let $E$ be a three-dimensional vector space and $S$ a linear application from $E$ to $E$ which associates to any vector $\vec{U}$ of $E$ its image $\vec{U}' = S(\vec{U})$. We are interested in the case where $S$ is an antisymmetric application defined by: $S(\vec{U}) = L \cdot \vec{U}$.

Find the matrix $L$ associated with the antisymmetric application $S$ by taking $s_1$, $s_2$, and $s_3$ as components of the vector $\vec{S}$.

\begin{correctionbox}

\end{correctionbox}

% ----------- EXERCISE 4 -------------
\section{}

To every point $P(x, y, z)$ of affine space, we associate the family of vector fields $\vec{V}_t(P)$ defined by: 
$$\vec{V}_t(P) = (3y - tz + 1, -3x + 2tz, tx + t^2y - 4/3); \quad t \in \mathbb{R}.$$

\begin{enumerate}
    \item Show that the only equiprojective fields are obtained for $t = 0$ and $t = 2$.
    \item Determine for each case the associated torsors by their reduction elements at point $O(0, 0, 0)$.
\end{enumerate}

\begin{correctionbox}

\end{correctionbox}

% ----------- EXERCISE 5 -------------
\section{}

We are given two sliders $(A, \vec{U})$ and $(B, \vec{V})$ such that: $A(1, 1, \alpha)$, $B(0, 2, 0)$, $\vec{U}(0, 0, \alpha)$ and $\vec{V}(\beta, 3, 0)$ where $\alpha$ and $\beta$ are real numbers. Let the torsor $T = (A, \vec{U}) + (B, \vec{V})$.

\begin{enumerate}
    \item Give the reduction elements of $T$ at point $O$.
    \item What is the necessary and sufficient condition for $T$ to be a slider?
    \item Determine the support of $T$.
\end{enumerate}

\begin{correctionbox}

\end{correctionbox}

% ----------- EXERCISE 6 -------------
\section{}

Let three sliders $(A_i, \vec{V}_i)$; $i = 1, 2, 3$ such that: $A_1(0, 0, a)$, $A_2(a, 0, 0)$ and $A_3(0, a, 0)$; $\vec{V}_1(2a, 0, 0)$, $\vec{V}_2(0, 3a, 0)$ and $\vec{V}_3(0, 0, \lambda a)$. We consider the torsor $L = \sum_{i=1}^{3} (A_i, \vec{V}_i)$. Determine the reduction elements of $L$ at an arbitrary point $P$ as well as its central axis.

\begin{correctionbox}

\end{correctionbox}

% ----------- EXERCISE 7 -------------
\section{}

Consider, in a Cartesian frame $R(O, X, Y, Z)$ with basis $(\vec{i}, \vec{j}, \vec{k})$, the following system of three bound vectors defined by:
\begin{align*}
A_1(1, 0, 0) &\quad \vec{V}_1 = a\vec{i} + b\vec{j} + \vec{k} \\
A_2(0, 1, 0) &\quad \vec{V}_2 = 2b\vec{i} - 2a\vec{j} - 3\vec{k} \\
A_3(0, 0, 1) &\quad \vec{V}_3 = -8\vec{i} + \vec{j} + c\vec{k}
\end{align*}

where $a$, $b$, and $c$ are real constants.

For what values of $a$, $b$, and $c$ is the system of three bound vectors equivalent to a couple? Determine then its moment.

\begin{correctionbox}

\end{correctionbox}

% ----------- EXERCISE 8 -------------
\section{}

We consider the vector field $\vec{W}$ which associates to every point $M(x, y, z)$ the vector $\vec{W}(M)$ with components:
\begin{align*}
W_x &= -1 - y - 2az \\
W_y &= 3 + x - 2az \\
W_z &= -2 + 2ax + a^2y
\end{align*}

where $a$ is a real parameter.

\begin{enumerate}
    \item For what values of $a$ is the field $\vec{W}$ an antisymmetric field?
    \item Determine, in each case, the reduction elements of the associated torsor.
\end{enumerate}

\begin{correctionbox}

\end{correctionbox}

\end{document}